\documentclass[11pt]{article}
\usepackage[margin=1in]{geometry}
\usepackage{amsmath,amssymb,mathtools}
\usepackage{fontspec,unicode-math}
\setmainfont{Latin Modern Roman}
\setsansfont{Latin Modern Sans}
\setmathfont{Latin Modern Math}
\usepackage{xurl,hyperref}
\hypersetup{colorlinks=true,urlcolor=blue,linkcolor=blue}
\newcommand{\sref}[2]{\hyperlink{source:#1:#2}{[S#2]}}
\newcommand{\cataloguescope}[1]{\paragraph{Recorded mathematical question.}#1}
\begin{document}
\section{Linear support of harmonic measure in recurrent planar triangulations}\label{statement}

\subsection{Definitions and mathematical statement}

\paragraph{Shared notation.}$\mathbb N=\{1,2,\ldots\}$, and $\mathbb Z,\mathbb Q,\mathbb R,\mathbb C$ denote the integers, rationals, reals and complex numbers. Any other conventions are specified locally.


Let $G$ be an infinite connected bounded-degree recurrent planar triangulation rooted at $o$, circle-packed in the whole Euclidean plane. Graph balls use unit-edge distance. For a finite connected vertex domain $U\ni o$, write $\partial_VU$ for vertices outside $U$ adjacent to $U$, and define harmonic measure $\omega_U(v)=\mathbb P_o(X_{\tau_U}=v)$ on $\partial_VU$, where $\tau_U=\inf\{t:X_t\notin U\}$ for simple random walk. Each boundary vertex corresponds to one circle. A random rooted graph is unimodular if for every nonnegative measurable function $F$ of a graph with two distinguished vertices,
\[\mathbb E\sum_vF(G,o,v)=\mathbb E\sum_vF(G,v,o).\]
\paragraph{Statement recorded in the supplied source.}
Are there radii $r_j\to\infty$, finite domains $U_j\supset B_G(o,r_j)$, boundary sets $A_j\subset\partial_VU_j$, and $\varepsilon_j\downarrow0$ such that
\[|A_j|\le r_j^{1+\varepsilon_j},\qquad\omega_{U_j}(A_j)\ge1-\varepsilon_j?\]
If this fails for general recurrent triangulations, does it hold almost surely under unimodularity? In particular, does the corresponding harmonic-measure conclusion hold for the UIPT? \sref{statement}{2}
\subsection{Short English statement}
Can almost all exit probability be concentrated on a nearly linear number of boundary circles along arbitrarily large scales, at least for unimodular triangulations or the UIPT?
\subsection{Sources}
\begin{description}
\item[\hypertarget{source:statement:1}{S1}] ULAM.AI and the UnsolvedMath Contributors, \emph{UnsolvedMath: A Curated Collection of Open Mathematics Problems}, catalogue entry. CC BY 4.0. \url{https://huggingface.co/datasets/ulamai/UnsolvedMath}.
\item[\hypertarget{source:statement:2}{S2}] Itai Benjamini, \emph{Random Planar Metrics}; original author manuscript preserved in Arquivo.pt capture of 31 December 2020, Section 3.3, harmonic-measure question and following qualification, p. 8. \url{https://arquivo.pt/noFrame/replay/20201231041544id_/http://www.wisdom.weizmann.ac.il/~itai/randomplanar2.pdf}.
\end{description}
\label{end:statement}
\end{document}
